\subsection{Implementation Deviations}
\label{sec:improvements}

The v4-sound reference implementation diverges from the spec at
several points.  Each deviation below is one the panel has reviewed
and recommends as a spec-side change rather than a one-off
implementation choice.

\one
\para{Per-policy vault contracts.}
The spec describes a single CM precompile keyed by $(\asset, \owner)$.
The implementation deploys an EIP-1167 minimal-proxy clone per token
policy and pins each policy's escrow + state to its own contract
address.
Two policies over the same underlying asset (e.g. \texttt{USDC-Bahamas}
and \texttt{USDC-EU} with disjoint \CE\ committees) get physically
disjoint trie subtrees; a bug in one clone cannot drain another's
collateral.
Recommendation: spec should make the \emph{vault contract} the
authoritative entry for a $(\asset, \mathsf{policyId})$ pair and
treat the precompile as a stateless verifier the vault delegates
into.

\one
\para{Per-account pending-credit cap.}
The spec fixes the chunked-pending counter ceiling at $2^{16}$ on the
chain.
The implementation lifts the cap onto $\acct.\mathsf{maxPending}$
(default 16), settable at $\texttt{configureAccount}$ time.
This lets a high-throughput recipient (e.g. a merchant) raise its
own ceiling without renegotiating the token policy, while a casual
recipient keeps a tight cap that bounds the decoder window on a
phone.
Recommendation: spec should expose the cap as account-state, not
chain-state, and tie its allowed range to the BSGS window the
recovery procedure declares.

\one
\para{Two-key compliance entity profile.}
The spec gives each \CE\ one DH key $\ceek$ and one ECDSA key $\vk$
and registers them separately.
The implementation bundles them in a single $\mathsf{CEProfile}$
struct $(\encryptpk, \authorizepk)$ stored on the
\texttt{TokenPolicy}, not on the account.
Putting the profile on the token makes the \CE\ identity portable
across all users of that token, lets the issuer rotate \CE\ keys
without touching $\sim$$N$ account rows, and reduces the per-account
state footprint to the floor enum.
Recommendation: move \CE\ identity onto the token policy in the
spec's account table.

\one
\para{Wrapped \CE\ authorization digest.}
The spec has the \CE\ sign the \texttt{TransferAuth} EIP-712 struct
directly with the same digest the on-chain check recomputes.
The implementation signs a domain-tagged wrap:
$\ceauthdigest := \hash(\mathtag{v4\text{-}ce\text{-}authorize} \cat
\transfercalldigest)$.
The wrap lets an auditor verify their own signature offline against
just $\transfercalldigest$ (a 32-byte input) without reconstructing
the full EIP-712 typed-data domain, which simplifies the
audit-tooling story.
The security argument is unchanged: collision-resistance of $\hash$
plus a unique domain tag.

\one
\para{Chunked-available with seed-derived blinding.}
The spec maintains a single twisted-ElGamal $\cbal$ for the main
balance and mandates a re-encryption step after every
\texttt{rollover} because the owner no longer knows $\rcur$.
The implementation chunks the main balance into four 32-bit limbs,
each blinded by $r_i := \HKDF(\dk;\, \mathtag{av} \cat \asset \cat
\nonce \cat i)$.
Because the wallet derives $r_i$ from the seed, it always knows the
opening; \texttt{Rerand} becomes optional (for key rotation) rather
than mandatory after every rollover.
Recommendation: spec should add \emph{chunked available} as an
alternative balance representation and document the
seed-derived-blinding contract so \texttt{Rerand} can be relaxed
from required to opt-in.

\one
\para{Domain-bound transfer randomness.}
The spec already binds $\chainid \cat \asset \cat \cenonce \cat i$
into the HKDF info for $r_i$ (\S\ref{impl:hkdf}).
The implementation initially defaulted the prover-side $\chainid$ to
a constant and hardcoded the verifier-side context to zero --- a
bug that produced verify-side digest mismatches under module
re-initialization.
Recommendation: the spec is right; the implementation has to take
\texttt{chain\_id} from the runtime context (passed through the
precompile FFI) and the prover from the connected wallet's network,
not from a constant.

\one
\para{Recipient lookup via configure-time event.}
The spec assumes the sender knows the recipient's $\ek$ out of band.
The implementation emits
$\mathsf{AccountConfigured}(\token, \mathsf{mmAddress},
\mathsf{tinyvmOwner})$ from the vault on every successful
\texttt{configureAccount}, indexed by both addresses.
A sender can paste either form (MM EOA or the internal account
address) and a one-call event scan resolves it.
The panel considers this a UX patch rather than a protocol
addition; the long-term direction is to collapse the two addresses
into one (see open questions, \S\ref{sec:open}).

\one
\para{Effective CE mode pinned by token + both endpoints.}
The implementation computes the effective mode as the lattice join
$\cemode_{\rm tok} \lor \cemode_{\rm out, snd} \lor
\cemode_{\rm in, rcv}$ and feeds it into the call digest the owner
signs.
The spec's $\cemode$ is per-token only; per-account inbound/outbound
floors are not in the data model.
Recommendation: spec should add account-side floors so a
high-stake recipient can demand $\Authorize$ even when the token
floor is $\Capsule$, without rejecting incoming transfers from
default senders.

\one
\para{Open: collapse \emph{two-address} confidential identity.}
\label{sec:open}
The implementation derives a seed-internal account address (the
v4-sound ``TinyVM owner'') from the EIP-191 signature in addition
to the EOA address that paid for gas, and stores per-account state
under the former.
This created a UX failure (users paste their MM address; the chain
has no row for it) that the $\mathsf{AccountConfigured}$ event
patches but does not solve.
Three resolutions are on the table:
%
\begin{enumerate}[itemsep=.5pt]
    \item Collapse: bind the account address to $\texttt{msg.sender}$
      at configure time; drop the seed-derived owner.  Verifier
      enforces $\texttt{msg.sender}$ ``owns'' the chosen $\ek$ via
      the configure-time signature; subsequent ops sign with the
      same EOA key.
    \item Canonical privacy address: keep two addresses but
      enforce that the dapp \emph{always} surfaces the privacy
      form for sharing.  Implementation cost is low; UX risk
      remains because users still find the EOA via block explorers.
    \item Stealth addresses (ERC-5564): a fresh per-transfer
      recipient address derived from the recipient's view tag.
      Eliminates the question of ``what is Bob's address''.
\end{enumerate}
%
Panel recommends (1) for v5 unless the stealth-address design is
ready in time.
